% problem-000098 4 化学反应与结构综合分析
\section*{4. 化学反应与结构综合分析}
（图：）
α

\subsection*{4-1}
根据该曲线，计算赖氨酸的等电点pI以及 $\mathrm { p H } = 9 . 6 0 $ 时赖氨酸4种分布形式的平衡浓度。

\subsection*{4-2}
实验室现有 0.10 mol L−1 $\mathrm { N a O H }$ 溶液、 $0 . 1 0 \mathrm { m o l } \mathrm { L } ^ { - 1 }$ HCl溶液、草酸基准物质 $\mathrm { ( H _ { 2 } C _ { 2 } O _ { 4 } { \cdot } 2 H _ { 2 } O }$ , 126.07 gmol−1)、邻苯⼆甲酸氢钾基准物质 $( \mathrm { K H C _ { 8 } H _ { 4 } O _ { 4 } } , 2 0 4 . 2 2 \mathrm { g } \mathrm { m o l ^ { - 1 } } )$ 、三(羟甲基)氨基甲烷基准物质（简称Tris,$\mathrm { ( H O C H _ { 2 } ) _ { 3 } C N H _ { 2 } }$ $K _ { \mathrm { b } } _ { \mathrm { b } } _ { \mathrm { \ell } } = 1 . 1 5 { \times } 1 0 ^ { - 6 }$ $1 2 1 . 1 4 \mathrm { g } \mathrm { m o l } ^ { - 1 } )$ ）、碳酸钠基准物质 $\mathrm { ( N a _ { 2 } C O _ { 3 } }$ , 105.99 g mol−1)，以及酚酞、甲基橙和甲基红等三种指⽰剂液。欲⽤酸碱滴定法测定某化学纯赖氨酸试样的含量。

\subsection*{4-2}
-1 指出合适的滴定剂并简要说明理由，写出滴定反应化学⽅程式。

\subsection*{4-2}
-2 计算滴定终点时溶液 $\mathrm { p H }$ ，根据计算结果选择最佳指⽰剂。

\subsection*{4-2}
-3 假设所选滴定剂的消耗量25 mL，通过计算所需基准物质的质量及其称量的相对误差确定最佳基准物质。

\subsection*{4-2}
-4 假设实验中准确称取� g赖氨酸试样（其摩尔质量⽤符号�表⽰），滴定⾄终点时消耗� mL浓度为� mol L−1的所选滴定剂标准溶液，写出该试样中赖氨酸质量分数(�)的计算式。
